\documentclass[11pt,a4paper]{article}

\usepackage{fontspec}
\usepackage{polyglossia}
\setmainlanguage{russian}
\setotherlanguage{english}

\IfFontExistsTF{Times New Roman}
{\setmainfont{Times New Roman}}
{\setmainfont{Libertinus Serif}}

\IfFontExistsTF{Menlo}
{\newfontfamily\cyrillicfonttt{Menlo}}
{\newfontfamily\cyrillicfonttt{DejaVu Sans Mono}}

\usepackage[a4paper,top=1.8cm,bottom=1.8cm,left=2.0cm,right=2.0cm]{geometry}
\usepackage{amsmath,amssymb}
\usepackage{booktabs}
\usepackage{array}
\usepackage{enumitem}
\usepackage{xcolor}
\usepackage{hyperref}
\usepackage{titlesec}

\hypersetup{colorlinks=true, linkcolor=black, urlcolor=blue}

\definecolor{hseblue}{HTML}{144F7B}
\definecolor{softblue}{HTML}{EAF4FA}
\definecolor{softgreen}{HTML}{EDF7F1}
\definecolor{softred}{HTML}{FBEDEA}
\definecolor{textgray}{HTML}{333333}

\setlength{\parindent}{0pt}
\setlength{\parskip}{0.52em}
\setlength{\emergencystretch}{3em}
\linespread{1.04}\selectfont

\titleformat{\section}
  {\Large\bfseries\color{hseblue}}
  {\thesection.}
  {0.5em}
  {}
\titleformat{\subsection}
  {\large\bfseries\color{hseblue}}
  {\thesubsection.}
  {0.5em}
  {}
\titleformat{\subsubsection}
  {\bfseries\color{textgray}}
  {}
  {0pt}
  {}

\setlist[itemize]{leftmargin=1.45em, topsep=2pt, itemsep=2pt}
\setlist[enumerate]{leftmargin=1.6em, topsep=2pt, itemsep=2pt}

\newcommand{\key}[1]{%
  \vspace{0.15em}
  \noindent\colorbox{softblue}{%
    \parbox{\dimexpr\linewidth-2\fboxsep\relax}{\textbf{Главная мысль.} #1}%
  }\vspace{0.2em}
}

\newcommand{\answer}[1]{%
  \vspace{0.15em}
  \noindent\colorbox{softgreen}{%
    \parbox{\dimexpr\linewidth-2\fboxsep\relax}{\textbf{Если спросят на защите.} #1}%
  }\vspace{0.2em}
}

\newcommand{\warn}[1]{%
  \vspace{0.15em}
  \noindent\colorbox{softred}{%
    \parbox{\dimexpr\linewidth-2\fboxsep\relax}{\textbf{Важно не перепутать.} #1}%
  }\vspace{0.2em}
}

\newcommand{\impl}[1]{%
  \vspace{0.15em}
  \noindent\colorbox{softblue}{%
    \parbox{\dimexpr\linewidth-2\fboxsep\relax}{\textbf{Как реализовано в коде.} #1}%
  }\vspace{0.2em}
}

\begin{document}

\begin{center}
{\Large\bfseries Подробный разбор методологии и практической части\par}
\vspace{0.3em}
{\large к курсовой работе «Исследование устойчивости моделей волатильности и ценообразования опционов в различных рыночных режимах»\par}
\vspace{0.5em}
Ковалёва Мария Сергеевна \\
НИУ ВШЭ, 2026
\end{center}

\vspace{0.8em}

Этот документ написан как подготовка к вопросам по методологии. В нём подробно разобраны: данные, модели ценообразования, оценки волатильности, метрики ошибок, TVNAE, премия раннего исполнения, статистические тесты, регрессия и проверки устойчивости. Формулы объясняются словами: что означает каждый символ, зачем применяется конкретный метод и как интерпретировать результат.

Обновление: в документ добавлены подробные пояснения к формулам, интуитивный вывод дифференциального уравнения Блэка--Шоулза, объяснение риск-нейтральной вероятности, обратной индукции и того, как расчёты реализованы в файлах \texttt{calculations.ipynb} и \texttt{options\_utils.py}.

\tableofcontents
\newpage

\section{Общая логика исследования}

\subsection{Что именно проверяется}

Исследование сравнивает \textbf{модельные цены} опционов с \textbf{рыночными расчётными ценами} Московской биржи. Для каждого наблюдения есть реальная биржевая цена опциона и несколько теоретических цен, рассчитанных по разным моделям. Затем строится ошибка:

\begin{equation}
e_i = \widehat{P}_i - P_i,
\end{equation}

где $\widehat{P}_i$ -- модельная цена, а $P_i$ -- рыночная расчётная цена. Если $e_i<0$, модель занизила цену относительно рынка. Если $e_i>0$, модель завысила цену.

Индекс $i$ означает конкретное наблюдение: один опцион в конкретную торговую дату. Знак «шляпка» в $\widehat{P}_i$ означает оценку, то есть цену, полученную из модели, а не напрямую наблюдаемую на бирже. Эта формула является базовой для всей практической части: все дальнейшие метрики и тесты строятся из $e_i$.

Главная идея работы не в том, чтобы доказать, что какая-то модель даёт «истинную» стоимость опциона. Рыночная расчётная цена сама может содержать ликвидностные эффекты, особенности биржевого расчёта и временные отклонения. Но она является наблюдаемым эмпирическим бенчмарком, поэтому по ней можно сравнивать устойчивость моделей.

\key{Проверяется не «абсолютная истинность» модели, а структура её ошибки: насколько ошибка велика, систематична и зависит от режима рынка, срока до экспирации и денежности.}

\subsection{Почему исследование построено именно так}

Рассматриваемые инструменты -- это опционы серии MM на фьючерс MXI. Фьючерс MXI является фьючерсом на Индекс МосБиржи, а опционы на него имеют американский стиль исполнения. Это определяет выбор моделей:

\begin{itemize}
    \item основная модель должна уметь учитывать досрочное исполнение, поэтому используется биномиальная модель Кокса--Росса--Рубинштейна;
    \item Black-76 применяется как европейский benchmark, потому что это стандартная аналитическая модель для опционов на фьючерсы, но она не учитывает early exercise;
    \item волатильность подаётся в модель тремя способами: HV\_21d, HV\_63d и GARCH;
    \item затем сравнивается, какая связка «модель цены + оценка волатильности» ближе к рыночной цене.
\end{itemize}

\answer{Если спросят, почему Black-76 вообще используется, если опционы американские: он нужен не как основная модель, а как контрольный европейский benchmark. Сравнение Black-76 и CRR показывает, насколько экономически важно право досрочного исполнения.}

\subsection{Практический pipeline расчётов}

Практическая часть устроена как последовательность шагов:

\begin{enumerate}
    \item Из исходной базы MOEX берутся дневные наблюдения по опционам и соответствующим фьючерсам.
    \item Фильтруются наблюдения с некорректными датами, неположительной ценой, неположительным страйком, нулевым сроком до экспирации и нераспознанным типом опциона.
    \item Для каждого наблюдения считаются $F$, $K$, $T$, DTE, moneyness и log-moneyness.
    \item По временным рядам фьючерсов считаются дневные лог-доходности.
    \item На их основе считаются HV\_21d, HV\_63d и GARCH-volatility.
    \item Для каждого наблюдения рассчитываются модельные цены: CRR + HV\_21d, CRR + HV\_63d, CRR + GARCH, Black-76 + HV\_63d.
    \item Рассчитываются ошибки и метрики качества.
    \item Результаты анализируются на полной выборке и в разрезах по vol regime, DTE и moneyness.
    \item Проводятся формальные проверки: тест средней ошибки, тесты Диболда--Мариано, OLS-регрессия ошибок и robustness checks.
\end{enumerate}

\section{Данные и переменные}

\subsection{Рабочая выборка}

Исходная база содержит 19 574 дневных наблюдения по опционам на фьючерс MXI за период с 3 января 2024 года по 19 мая 2026 года. После фильтрации остаётся 19 565 наблюдений. Общая сопоставимая выборка для сравнения трёх американских моделей содержит 18 058 наблюдений.

Сопоставимая выборка нужна, чтобы все модели сравнивались на одном и том же наборе наблюдений. Если сравнивать модели на разных строках, различия в ошибке могут возникнуть не из-за качества модели, а из-за разного состава данных.

\answer{Если спросят, почему наблюдений стало меньше: часть наблюдений теряется из-за требований к доступности всех оценок волатильности, особенно GARCH, которому нужно накопить минимум 63 предыдущих наблюдения.}

\subsection{Основные обозначения}

\begin{itemize}
    \item $F_t$ -- расчётная цена базового фьючерса в дату $t$.
    \item $K$ -- страйк опциона.
    \item $T$ -- время до экспирации в годах.
    \item DTE -- срок до экспирации в календарных днях.
    \item $r$ -- безрисковая ставка; в базовой спецификации используется 16\% годовых.
    \item $\sigma$ -- волатильность, подаваемая в модель ценообразования.
    \item $P_i$ -- рыночная расчётная цена опциона.
    \item $\widehat{P}_i$ -- модельная цена опциона.
\end{itemize}

\subsection{Денежность и логарифмическая денежность}

Денежность задаётся как отношение цены фьючерса к страйку:

\begin{equation}
\text{moneyness}_i = \frac{F_i}{K_i}.
\end{equation}

Здесь $F_i$ -- цена фьючерса, а $K_i$ -- страйк. Если отношение равно 1, фьючерсная цена примерно равна страйку, и опцион находится около денег. Если для call-опциона $F_i/K_i>1$, исполнение уже имеет положительный смысл: можно купить дешевле страйка относительно рынка фьючерса. Для put-опциона экономическая интерпретация обратная, но само отношение всё равно удобно как единая координата положения контракта.

Логарифмическая денежность:

\begin{equation}
\ln\left(\frac{F_i}{K_i}\right).
\end{equation}

Логарифм используется потому, что он симметризует относительные изменения. Например, переход от 100 к 110 и от 100 к 90 в обычных процентах не полностью симметричен, а логарифмические отношения лучше подходят для моделей, где цены описываются через лог-доходности.

Абсолютная логарифмическая денежность:

\begin{equation}
\left|\ln\left(\frac{F_i}{K_i}\right)\right|.
\end{equation}

Она показывает, насколько контракт далеко от состояния at-the-money независимо от направления. Если значение близко к нулю, опцион около денег. Если значение большое, опцион далеко от денег или глубоко в деньгах.

\warn{Для коллов и путов «в деньгах» определяется по-разному, но абсолютная log-moneyness в регрессии используется как расстояние от ATM-зоны, а не как знак прибыльности опциона.}

\section{Внутренняя и временна́я стоимость опциона}

\subsection{Внутренняя стоимость}

Внутренняя стоимость -- это то, что держатель получил бы при немедленном исполнении опциона.

Для call-опциона:

\begin{equation}
H_{\text{call}}(F,K)=\max(F-K,0).
\end{equation}

В этой формуле $F-K$ -- выгода от исполнения call прямо сейчас: держатель получает возможность купить по страйку $K$, когда фьючерс стоит $F$. Если $F<K$, исполнение невыгодно, поэтому внутренняя стоимость не становится отрицательной, а берётся ноль.

Для put-опциона:

\begin{equation}
H_{\text{put}}(F,K)=\max(K-F,0).
\end{equation}

Для put логика обратная: держатель выигрывает, когда может продать по страйку $K$, если рыночная фьючерсная цена $F$ ниже. Поэтому выгода равна $K-F$, но опять же не ниже нуля.

Если внутренняя стоимость положительна, опцион находится in-the-money. Если она равна нулю, опцион out-of-the-money или at-the-money.

\subsection{Временна́я стоимость}

Временна́я стоимость -- это часть цены сверх внутренней стоимости:

\begin{equation}
TV_i = P_i - H(F_i,K_i).
\end{equation}

Здесь $P_i$ -- фактическая рыночная цена опциона, а $H(F_i,K_i)$ -- его внутренняя стоимость. Если опцион стоит 500 пунктов, а немедленное исполнение даёт 430 пунктов, то 70 пунктов -- это временна́я стоимость. Она платится за шанс, что до экспирации ситуация улучшится.

Экономически временна́я стоимость отражает ценность неопределённости до экспирации. Даже если опцион сейчас невыгодно исполнять, до экспирации цена фьючерса может измениться, и эта возможность имеет стоимость.

\key{Внутренняя стоимость определяется текущим положением $F$ и $K$, а временна́я стоимость зависит от срока, волатильности, ставки и ожиданий рынка. Поэтому качество волатильностной модели лучше оценивать именно через временну́ю стоимость.}

\section{Оценка волатильности}

\subsection{Почему волатильность центральна}

В опционных моделях волатильность определяет ширину распределения будущей цены базового актива. Чем выше волатильность, тем выше вероятность сильного движения цены. Для держателя опциона это ценно, потому что убыток ограничен премией, а потенциальная выгода может расти.

Поэтому при прочих равных рост $\sigma$ обычно увеличивает цену и call, и put. Главная эмпирическая проблема: будущая волатильность неизвестна, её нужно оценить.

\subsection{Дневные лог-доходности}

Для каждого базового фьючерса строится временной ряд дневных логарифмических доходностей:

\begin{equation}
r_t = \ln\left(\frac{F_t}{F_{t-1}}\right).
\end{equation}

В этой формуле $F_t$ -- цена фьючерса сегодня, $F_{t-1}$ -- цена вчера. Если цена выросла, отношение больше 1 и лог-доходность положительная. Если цена упала, отношение меньше 1 и лог-доходность отрицательная. Лог-доходность приблизительно равна обычной процентной доходности при малых изменениях, но удобнее математически.

Лог-доходности удобны, потому что они аддитивны во времени и естественно связаны с моделями геометрического броуновского движения.

\subsection{Историческая волатильность}

Историческая волатильность рассчитывается как стандартное отклонение прошлых доходностей на скользящем окне:

\begin{equation}
\hat{\sigma}^2_t =
\frac{1}{n-1}
\sum_{j=0}^{n-1}
\left(r_{t-j}-\bar r_t\right)^2,
\qquad
\bar r_t =
\frac{1}{n}
\sum_{j=0}^{n-1} r_{t-j}.
\end{equation}

Здесь $n$ -- длина окна, например 21 или 63 дня. Величина $\bar r_t$ -- средняя доходность внутри этого окна. Выражение $(r_{t-j}-\bar r_t)^2$ показывает, насколько конкретная доходность отклоняется от среднего. Сумма таких квадратов измеряет разброс доходностей, а деление на $n-1$ даёт выборочную оценку дисперсии. Корень из дисперсии -- это стандартное отклонение, то есть дневная волатильность.

Дневная волатильность годируется:

\begin{equation}
HV_{n,t}=\sqrt{252}\,\hat{\sigma}_t.
\end{equation}

В этой формуле $HV_{n,t}$ -- годовая историческая волатильность на дату $t$, рассчитанная по окну длиной $n$. Если дневная волатильность равна 1\%, то годовая при стандартном масштабировании будет примерно $1\%\cdot\sqrt{252}\approx15.9\%$.

Множитель $\sqrt{252}$ используется потому, что дисперсия масштабируется линейно по времени, а стандартное отклонение масштабируется как корень из времени. 252 -- стандартное приближение числа торговых дней в году.

\subsection{HV\_21d}

В работе:

\begin{equation}
\sigma_{21,t} =
\sqrt{252}\,\operatorname{sd}(r_{t-20},\ldots,r_t).
\end{equation}

Здесь $\operatorname{sd}$ -- стандартное отклонение 21 последней дневной лог-доходности: от $r_{t-20}$ до $r_t$. Индекс 21 означает не календарные дни, а торговые наблюдения.

Это короткое окно примерно на один торговый месяц. Его смысл -- быстро реагировать на недавние изменения рынка. Если волатильность резко выросла, HV\_21d быстрее подхватит этот рост.

Минус: короткое окно шумное. Одна-две экстремальные доходности могут сильно изменить оценку. Для длинных опционов такая краткосрочная оценка может плохо соответствовать горизонту контракта.

\subsection{HV\_63d}

В работе:

\begin{equation}
\sigma_{63,t} =
\sqrt{252}\,\operatorname{sd}(r_{t-62},\ldots,r_t).
\end{equation}

Здесь используется уже 63 последние дневные доходности. Формула такая же, как для HV\_21d, но окно длиннее, поэтому отдельный выброс влияет на итоговую волатильность слабее.

Это окно примерно на один квартал. Оно более сглаженное и устойчивое. В данной выборке много длинных опционов, поэтому HV\_63d оказывается сильным baseline: она меньше реагирует на краткосрочный шум и лучше отражает более длинный горизонт.

\answer{Если спросят, почему выбраны 21 и 63 дня: 21 торговый день примерно соответствует месяцу, 63 -- кварталу. Это стандартные практические горизонты для краткосрочной и более устойчивой исторической волатильности.}

\subsection{GARCH(1,1)}

\subsubsection{Идея модели}

GARCH нужен, потому что финансовая волатильность не постоянна. На рынках часто наблюдается кластеризация волатильности: большие движения следуют за большими, спокойные периоды -- за спокойными. Доходности сами могут иметь слабую автокорреляцию, но квадраты доходностей и абсолютные доходности часто зависимы во времени.

\subsubsection{Уравнения GARCH}

Базовая модель:

\begin{equation}
r_t = \mu + \varepsilon_t,
\end{equation}

Здесь $r_t$ -- доходность в момент $t$, $\mu$ -- условное среднее доходности, а $\varepsilon_t$ -- неожиданная часть доходности, то есть шок. Если фактическая доходность выше ожидаемой, шок положительный; если ниже -- отрицательный.

\begin{equation}
\varepsilon_t = \sigma_t z_t,
\end{equation}

В этой записи шок раскладывается на масштаб и случайную стандартизованную величину. $\sigma_t$ -- условная волатильность в момент $t$, а $z_t$ -- случайная величина со средним 0 и дисперсией 1. То есть $z_t$ отвечает за направление и стандартизированный размер шока, а $\sigma_t$ масштабирует его в реальные рыночные единицы.

\begin{equation}
\sigma_t^2 =
\omega
+ \alpha \varepsilon_{t-1}^2
+ \beta \sigma_{t-1}^2.
\end{equation}

Это главное уравнение GARCH. Оно говорит, что сегодняшняя условная дисперсия зависит от трёх частей: постоянной базы $\omega$, вчерашнего шока $\varepsilon_{t-1}^2$ и вчерашней условной дисперсии $\sigma_{t-1}^2$. Квадрат шока используется потому, что важен размер неожиданного движения, а не его знак: и сильный рост, и сильное падение повышают оценку волатильности.

В практической реализации работы используется нулевое условное среднее:

\begin{equation}
r_t = \varepsilon_t,
\end{equation}

Это означает, что $\mu$ в предыдущей формуле приравнивается к нулю. Для дневных доходностей это распространённое упрощение: средняя дневная доходность мала по сравнению с дневной волатильностью, а для ценообразования здесь важнее динамика дисперсии.

потому что для дневных фьючерсных доходностей среднее обычно мало относительно волатильности, а задача состоит именно в прогнозе дисперсии.

\subsubsection{Смысл параметров}

\begin{itemize}
    \item $\omega$ -- базовый уровень дисперсии. Если бы не было новых шоков, он поддерживает положительную условную дисперсию.
    \item $\alpha$ -- реакция на новый шок. Чем больше $\alpha$, тем сильнее вчерашний большой остаток повышает сегодняшнюю волатильность.
    \item $\beta$ -- инерция прошлой волатильности. Чем больше $\beta$, тем медленнее волатильность возвращается к нормальному уровню.
    \item $\alpha+\beta$ -- персистентность. Если сумма близка к единице, шоки долго сохраняются.
\end{itemize}

Для стационарности обычно требуется:

\begin{equation}
\alpha+\beta<1.
\end{equation}

Это условие означает, что влияние шоков со временем затухает. Если $\alpha+\beta$ было бы равно или больше 1, волатильность могла бы не возвращаться к долгосрочному уровню, и модель стала бы плохо интерпретируемой для устойчивого прогноза.

Безусловная долгосрочная дисперсия:

\begin{equation}
\bar{\sigma}^2 =
\frac{\omega}{1-\alpha-\beta}.
\end{equation}

Здесь $\bar{\sigma}^2$ -- уровень дисперсии, к которому модель стремится в долгосрочном равновесии. Чем ближе $\alpha+\beta$ к 1, тем больше знаменатель приближается к нулю и тем медленнее модель возвращается к долгосрочному уровню после шока.

\subsubsection{Прогноз дисперсии}

Одношаговый прогноз:

\begin{equation}
\hat{\sigma}_{t+1|t}^2
=
\omega
+ \alpha\varepsilon_t^2
+ \beta\sigma_t^2.
\end{equation}

Индекс $t+1|t$ читается как «прогноз на следующий период, сделанный в момент $t$». В формуле используются только данные, известные к моменту $t$: текущий шок $\varepsilon_t$ и текущая условная дисперсия $\sigma_t^2$.

Многошаговая форма:

\begin{equation}
\hat{\sigma}^2_{t+h|t}
=
\bar{\sigma}^2
+(\alpha+\beta)^{h-1}
\left(\sigma_t^2-\bar{\sigma}^2\right).
\end{equation}

Здесь $h$ -- горизонт прогноза. Формула показывает mean reversion: прогноз будущей дисперсии равен долгосрочному уровню плюс отклонение текущей дисперсии от долгосрочной, умноженное на затухающий коэффициент $(\alpha+\beta)^{h-1}$.

Для ценообразования в работе используется однодневный прогноз, приведённый к годовому масштабу:

\begin{equation}
\hat{\sigma}_{\text{ann},t}
=
\sqrt{252}\,\hat{\sigma}_{t+1|t}.
\end{equation}

Здесь $\hat{\sigma}_{\text{ann},t}$ -- годовая волатильность, которую затем можно подать в CRR. Модели ценообразования используют годовую $\sigma$, потому что срок $T$ тоже выражен в годах.

\subsubsection{Как GARCH реализован практически}

В коде GARCH оценивается отдельно для каждого базового фьючерса:

\begin{itemize}
    \item используется минимум 63 наблюдения перед первым прогнозом;
    \item окно расширяющееся, то есть на дату $t$ модель использует только прошлые данные;
    \item параметры переоцениваются каждые 21 торговый день;
    \item между переоценками дисперсия обновляется рекурсивно;
    \item доходности умножаются на 100 при оценивании, а затем прогноз волатильности переводится обратно в доли и годируется.
\end{itemize}

\key{Главная защита от look-ahead bias: прогноз на дату $t$ строится только по информации, доступной до этой даты. Будущие доходности не используются.}

\impl{В \texttt{options\_utils.py} GARCH не написан полностью вручную: параметры модели оцениваются стандартной библиотекой \texttt{arch} через \texttt{arch\_model(..., mean="Zero", vol="GARCH", p=1, q=1, dist="normal")}. Но логика expanding window, переоценка каждые 21 день, рекурсивное обновление прогнозной дисперсии и перевод в годовую волатильность написаны вручную в функции \texttt{compute\_expanding\_garch\_forecast}.}

\answer{Если спросят, почему GARCH не победил всегда: он даёт краткосрочный условный прогноз. Он полезен для коротких DTE, но длинный опцион зависит от ожидаемой волатильности на месяцы вперёд, поэтому однодневный прогноз не всегда оптимален.}

\section{Модели ценообразования}

\subsection{Black--Scholes как теоретическая база}

Модель Блэка--Шоулза строится для ответа на вопрос: сколько должен стоить европейский опцион сегодня, если известны текущая цена базового актива, страйк, срок до экспирации, безрисковая ставка и волатильность. Модель выдаёт теоретическую цену европейского call или put. Дополнительно из неё можно получить чувствительности опциона -- например, дельту, гамму, вегу, тету и ро.

Смысл модели: опцион можно оценить не через субъективный прогноз «вырастет рынок или упадёт», а через отсутствие арбитража. Если можно динамически собрать портфель из базового актива и безрискового счёта, который повторяет выплату опциона, то цена опциона должна совпадать со стоимостью такого реплицирующего портфеля. Иначе возникла бы арбитражная возможность.

\subsubsection{Предпосылка о движении цены}

Модель предполагает, что цена базового актива следует геометрическому броуновскому движению:

\begin{equation}
dS_t = \mu S_t dt + \sigma S_t dW_t,
\end{equation}

где $S_t$ -- цена актива, $\mu$ -- ожидаемая доходность, $\sigma$ -- постоянная волатильность, $W_t$ -- броуновское движение.

Здесь $dS_t$ -- малое изменение цены за малый промежуток времени $dt$. Часть $\mu S_tdt$ -- ожидаемое плавное изменение цены, или drift. Часть $\sigma S_tdW_t$ -- случайный шок. Чем больше $\sigma$, тем сильнее случайные колебания. Броуновское движение $W_t$ -- математическая модель непрерывного случайного шума.

Почему это называется геометрическим броуновским движением: случайной является относительная доходность, а не абсолютное изменение цены. Поэтому при такой модели цена остаётся положительной, а логарифм цены имеет нормальное распределение.

\subsubsection{Что означает цена опциона как функция}

Цена опциона обозначается как:

\begin{equation}
V=V(S,t).
\end{equation}

Это означает, что цена опциона зависит от текущей цены базового актива $S$ и текущего времени $t$. Остальные параметры -- страйк $K$, ставка $r$, волатильность $\sigma$ и срок экспирации -- фиксированы внутри задачи. На экспирации условие известно: для call это $\max(S_T-K,0)$, для put -- $\max(K-S_T,0)$.

\subsubsection{Самый простой вывод дифференциального уравнения}

Шаг 1. Записываем изменение цены опциона. Поскольку $V$ зависит от случайной цены $S_t$ и времени, используется лемма Ито:

\begin{equation}
dV =
\frac{\partial V}{\partial t}dt
+
\frac{\partial V}{\partial S}dS
+
\frac{1}{2}
\frac{\partial^2 V}{\partial S^2}
(dS)^2.
\end{equation}

Здесь $\frac{\partial V}{\partial t}$ показывает чувствительность цены опциона ко времени, $\frac{\partial V}{\partial S}$ -- чувствительность к цене базового актива, а $\frac{\partial^2V}{\partial S^2}$ -- кривизну зависимости цены опциона от базового актива. В финансовой терминологии это тета, дельта и гамма, хотя в самой формуле мы используем математические производные.

Шаг 2. Подставляем динамику $dS_t=\mu S_tdt+\sigma S_tdW_t$. В стохастическом исчислении используется правило $(dW_t)^2=dt$, поэтому:

\begin{equation}
(dS_t)^2 = \sigma^2S_t^2dt.
\end{equation}

После подстановки получаем:

\begin{equation}
dV =
\left(
\frac{\partial V}{\partial t}
+\mu S\frac{\partial V}{\partial S}
+\frac{1}{2}\sigma^2S^2
\frac{\partial^2V}{\partial S^2}
\right)dt
+
\sigma S\frac{\partial V}{\partial S}dW.
\end{equation}

Эта формула говорит: изменение цены опциона состоит из ожидаемой части перед $dt$ и случайной части перед $dW$.

Шаг 3. Собираем портфель из одного опциона и короткой позиции в $\Delta$ единицах базового актива:

\begin{equation}
\Pi = V - \Delta S,
\qquad
\Delta = \frac{\partial V}{\partial S}.
\end{equation}

Здесь $\Pi$ -- стоимость портфеля. Мы берём $\Delta$ равной дельте опциона. Это не случайный трюк: случайная часть в $dV$ равна $\sigma S \frac{\partial V}{\partial S}dW$, а случайная часть в $\Delta dS$ равна $\Delta\sigma SdW$. Если $\Delta=\frac{\partial V}{\partial S}$, эти случайные части взаимно уничтожаются.

Шаг 4. Считаем изменение портфеля:

\begin{equation}
d\Pi=dV-\Delta dS.
\end{equation}

После сокращения случайной части остаётся:

\begin{equation}
d\Pi =
\left(
\frac{\partial V}{\partial t}
+\frac{1}{2}\sigma^2S^2
\frac{\partial^2V}{\partial S^2}
\right)dt.
\end{equation}

Здесь важно: из формулы исчезла $\mu$, то есть ожидаемая доходность актива. Это один из центральных результатов Блэка--Шоулза: цена опциона в безарбитражной модели не зависит от субъективной ожидаемой доходности базового актива.

Шаг 5. Раз портфель безрисковый, он должен зарабатывать безрисковую ставку:

\begin{equation}
d\Pi = r\Pi dt = r\left(V-\frac{\partial V}{\partial S}S\right)dt.
\end{equation}

Если бы безрисковый портфель зарабатывал больше или меньше $r$, можно было бы получить арбитраж: занять деньги под ставку $r$ и вложить в портфель с большей доходностью либо наоборот.

Шаг 6. Приравниваем две записи для $d\Pi$:

\begin{equation}
\frac{\partial V}{\partial t}
+\frac{1}{2}\sigma^2S^2
\frac{\partial^2V}{\partial S^2}
=
rV-rS\frac{\partial V}{\partial S}.
\end{equation}

Переносим всё в левую часть и получаем дифференциальное уравнение Блэка--Шоулза:

\begin{equation}
\frac{\partial V}{\partial t}
+
\frac{1}{2}\sigma^2 S^2
\frac{\partial^2 V}{\partial S^2}
+
rS\frac{\partial V}{\partial S}
-rV=0.
\end{equation}

Это уравнение говорит: если цена опциона не удовлетворяет такому соотношению между временем, дельтой, гаммой, ставкой и волатильностью, то можно построить арбитражный портфель.

\key{Самая важная идея вывода: мы не угадываем будущую доходность актива. Мы строим портфель, где случайный риск локально убран через дельта-хеджирование, и поэтому этот портфель должен приносить безрисковую ставку.}

\subsubsection{Решение для европейских опционов}

Для европейского call:

\begin{equation}
C=S_0N(d_1)-Ke^{-rT}N(d_2),
\end{equation}

Здесь $C$ -- цена call сегодня, $S_0$ -- текущая цена базового актива, $K$ -- страйк, $e^{-rT}$ -- дисконтирующий множитель, а $N(\cdot)$ -- функция распределения стандартной нормальной случайной величины.

\begin{equation}
d_1 =
\frac{\ln(S_0/K)+(r+\sigma^2/2)T}{\sigma\sqrt{T}},
\qquad
d_2=d_1-\sigma\sqrt{T}.
\end{equation}

Величины $d_1$ и $d_2$ -- это нормированные расстояния в логнормальной модели цены. Очень грубо: они показывают, насколько текущая цена, страйк, ставка, срок и волатильность вместе делают опцион вероятно исполняемым. $N(d_2)$ часто интерпретируют как риск-нейтральную вероятность того, что call окажется in-the-money. $N(d_1)$ связан с дельтой call и весом базового актива в реплицирующем портфеле.

Для европейского put:

\begin{equation}
P=Ke^{-rT}N(-d_2)-S_0N(-d_1).
\end{equation}

Здесь $P$ -- цена европейского put. Формула похожа на call, но логика выплаты обратная: put приносит прибыль, когда цена базового актива ниже страйка.

\subsubsection{Вероятностный смысл}

Формулу Блэка--Шоулза можно понимать через риск-нейтральное ожидание:

\begin{equation}
V_0=e^{-rT}\mathbb{E}^{Q}\left[\text{payoff at }T\right].
\end{equation}

Здесь $\mathbb{E}^{Q}$ -- ожидание не под реальными вероятностями, а под риск-нейтральной мерой $Q$. Это специальный набор вероятностей, при котором дисконтированные цены активов ведут себя как мартингалы, то есть не имеют «бесплатного» ожидаемого роста сверх безрисковой ставки.

Для call:

\begin{equation}
C=e^{-rT}\mathbb{E}^{Q}\left[(S_T-K)^+\right].
\end{equation}

Эта запись означает: берём будущую выплату call, усредняем её по риск-нейтральным вероятностям и дисконтируем назад на сегодня.

В самой практической части эта модель не является основной, потому что исследуются опционы на фьючерс и американский стиль исполнения. Но она важна как теоретическая основа Black-76 и вообще классического опционного ценообразования.

\warn{Black--Scholes в чистом виде предназначен для европейских опционов на спотовый актив. В работе для фьючерсных опционов используется его адаптация -- Black-76.}

\subsection{Black-76}

\subsubsection{Почему Black-76 подходит для фьючерсов}

В Black-76 базовой переменной является не спотовая цена $S_t$, а фьючерсная цена $F_t$. В риск-нейтральной мере фьючерсная цена не имеет дрейфа:

\begin{equation}
dF_t = \sigma F_t dW_t.
\end{equation}

Здесь $F_t$ -- фьючерсная цена, $dF_t$ -- её малое изменение, $\sigma$ -- волатильность фьючерсной цены, $dW_t$ -- случайный шок. В формуле нет члена $\mu F_tdt$, потому что под риск-нейтральной мерой фьючерсная цена не имеет ожидаемого drift.

Интуитивно: фьючерс не требует первоначальной инвестиции в базовый актив, поэтому под риск-нейтральной мерой его ожидаемое изменение после учёта маржинальной механики равно нулю.

\subsubsection{Формулы Black-76}

Для европейского call на фьючерс:

\begin{equation}
C =
e^{-rT}
\left[
F_0N(d_1)-KN(d_2)
\right].
\end{equation}

Здесь $C$ -- цена европейского call на фьючерс, $F_0$ -- текущая фьючерсная цена, $K$ -- страйк, $T$ -- срок до экспирации в годах, $r$ -- безрисковая ставка, $e^{-rT}$ -- дисконтирование будущей ожидаемой выплаты к текущему моменту.

Для европейского put:

\begin{equation}
P =
e^{-rT}
\left[
KN(-d_2)-F_0N(-d_1)
\right].
\end{equation}

Здесь $P$ -- цена европейского put. Внутри квадратных скобок стоит ожидаемая риск-нейтральная выплата, а множитель $e^{-rT}$ переводит её в сегодняшнюю стоимость.

Здесь:

\begin{equation}
d_1 =
\frac{\ln(F_0/K)+\frac{1}{2}\sigma^2T}
{\sigma\sqrt{T}},
\qquad
d_2=d_1-\sigma\sqrt{T}.
\end{equation}

В этой формуле $\ln(F_0/K)$ показывает относительное положение фьючерса и страйка, $\frac{1}{2}\sigma^2T$ корректирует логнормальное распределение, а знаменатель $\sigma\sqrt{T}$ нормирует расстояние на масштаб неопределённости до экспирации. Чем выше волатильность или срок, тем больше неопределённость и тем меньше «жёсткость» расстояния между $F_0$ и $K$.

\subsubsection{Смысл показателей d1 и d2}

Строго $d_1$ и $d_2$ -- нормированные величины в риск-нейтральном логнормальном распределении. Интуитивно:

\begin{itemize}
    \item $d_2$ связан с риск-нейтральной вероятностью оказаться in-the-money к экспирации;
    \item $d_1$ связан с весом фьючерсной компоненты в ожидаемой выплате;
    \item $N(d)$ -- вероятность или вес под стандартным нормальным распределением.
\end{itemize}

\subsubsection{Ограничение Black-76 в работе}

Black-76 считает только европейский опцион. В нём нет шага, где держатель сравнивает «исполнить сейчас» и «ждать дальше». Поэтому модель может дать цену ниже внутренней стоимости. Для американского опциона это недопустимо, потому что держатель может немедленно исполнить опцион.

\key{Black-76 полезен как benchmark без досрочного исполнения. Если Black-76 сильно хуже CRR, значит early exercise действительно имеет экономическую ценность.}

\subsection{Биномиальная модель Кокса--Росса--Рубинштейна}

\subsubsection{Для чего строится CRR}

Модель Кокса--Росса--Рубинштейна строится для того же общего вопроса, что и Black--Scholes: сколько должен стоить опцион сегодня. Но CRR решает задачу численно, а не аналитической формулой. Это особенно полезно, когда у опциона есть особенности, которые трудно встроить в закрытую формулу, например досрочное исполнение.

Модель выдаёт цену опциона в начальный момент. Если сохранить всё дерево, можно также увидеть, в каких узлах рационально держать опцион дальше, а в каких выгодно исполнить его досрочно.

\subsubsection{Идея дерева}

CRR -- численная модель. Время до экспирации делится на $N$ маленьких шагов:

\begin{equation}
\Delta t = \frac{T}{N}.
\end{equation}

Здесь $T$ -- общий срок до экспирации в годах, $N$ -- число шагов дерева, а $\Delta t$ -- длина одного маленького шага. Чем больше $N$, тем мельче сетка времени и тем ближе дерево к непрерывной модели цены.

В работе используется $N=100$ шагов. На каждом шаге фьючерсная цена может подняться в $u$ раз или упасть в $d$ раз:

\begin{equation}
u = e^{\sigma\sqrt{\Delta t}},
\qquad
d = \frac{1}{u}.
\end{equation}

Здесь $u$ -- множитель движения вверх, $d$ -- множитель движения вниз, $\sigma$ -- годовая волатильность, а $\sqrt{\Delta t}$ переводит годовую волатильность в масштаб одного шага. Чем выше волатильность или длиннее шаг, тем дальше друг от друга возможные цены после движения вверх и вниз.

Такой выбор делает дерево рекомбинирующим: движение вверх, потом вниз приводит к той же цене, что движение вниз, потом вверх. Например, $Fud=Fdu=F$, потому что $d=1/u$. Это резко снижает вычислительную сложность: не нужно хранить все возможные пути отдельно, достаточно хранить узлы дерева.

\subsubsection{Почему биномиальное дерево -- нормальная аппроксимация}

На одном шаге модель очень грубая: цена может только вырасти или упасть. Но если шагов много, сумма большого числа маленьких независимых движений начинает вести себя похоже на нормальное распределение лог-доходности. Это интуиция центральной предельной теоремы.

В CRR лог-цена после многих шагов складывается из большого количества малых приращений: движение вверх даёт $+\sigma\sqrt{\Delta t}$ в логах, движение вниз даёт $-\sigma\sqrt{\Delta t}$. Когда $N$ растёт, дискретная биномиальная модель всё лучше приближает непрерывное геометрическое броуновское движение, лежащее в основе Black--Scholes.

\key{CRR логична как «дискретная версия» непрерывной модели: мы заменяем бесконечно много возможных движений цены большим числом маленьких шагов вверх/вниз. Чем мельче шаг, тем лучше аппроксимация.}

\subsubsection{Что такое риск-нейтральная вероятность}

Риск-нейтральная вероятность -- это не реальная вероятность роста цены. Это искусственная вероятность, которая используется для безарбитражного ценообразования. Она подбирается так, чтобы ожидаемая доходность базового актива в модели соответствовала безрисковой ставке или, для фьючерса, чтобы ожидаемая фьючерсная цена была мартингалом.

Простой пример для акции. Пусть акция сегодня стоит 100. Через один шаг она может стоить 120 или 90. Безрисковый рост за шаг равен 5\%, то есть 100 превращается в 105. Риск-нейтральная вероятность $q$ выбирается из условия:

\begin{equation}
105=q\cdot120+(1-q)\cdot90.
\end{equation}

Отсюда:

\begin{equation}
q=\frac{105-90}{120-90}=0.5.
\end{equation}

Это не значит, что реальная вероятность роста акции равна 50\%. Реальная вероятность может быть 60\%, 40\% или другой. Но для цены производного инструмента в безарбитражной модели мы используем такие вероятности, при которых базовый актив в среднем растёт по безрисковой ставке. Так цена опциона становится совместимой с ценой базового актива и безрисковой ставкой.

\answer{Если спросят «почему можно использовать не реальные вероятности»: потому что цена опциона в безарбитражной теории определяется репликацией и отсутствием арбитража, а не субъективным прогнозом инвестора. Риск-нейтральные вероятности -- удобный математический способ получить ту же цену, что и реплицирующий портфель.}

\subsubsection{Риск-нейтральная вероятность для фьючерса}

Для обычного актива в CRR часто используется:

\begin{equation}
q=\frac{e^{r\Delta t}-d}{u-d}.
\end{equation}

Здесь $e^{r\Delta t}$ -- безрисковый рост за один шаг, $u$ -- движение вверх, $d$ -- движение вниз. Формула выбирает $q$ так, чтобы ожидаемый рост базового актива под риск-нейтральной вероятностью был равен безрисковому росту. Это ровно та же идея, что в примере выше с 100, 120 и 90.

Но для фьючерсного опциона применяется условие мартингальности фьючерсной цены. В риск-нейтральной мере ожидаемая цена фьючерса на следующем шаге равна текущей:

\begin{equation}
F = qFu + (1-q)Fd.
\end{equation}

Левая часть $F$ -- текущая фьючерсная цена. Правая часть -- ожидаемая фьючерсная цена на следующем шаге: с вероятностью $q$ цена становится $Fu$, с вероятностью $1-q$ становится $Fd$. Для фьючерса под риск-нейтральной мерой это ожидание должно равняться текущей цене.

Делим на $F$:

\begin{equation}
1 = qu+(1-q)d.
\end{equation}

После деления остаются только относительные множители вверх и вниз. Это показывает, что сама величина текущей цены $F$ не влияет на риск-нейтральную вероятность: важны только пропорции движения.

Отсюда:

\begin{equation}
q=\frac{1-d}{u-d}.
\end{equation}

Здесь ставка $r$ исчезает из числителя, потому что фьючерсная цена под риск-нейтральной мерой не должна иметь drift. Но ставка не исчезает из всей модели: будущие значения опциона всё равно дисконтируются через $e^{-r\Delta t}$.

\warn{Ставка $r$ не входит в риск-нейтральную вероятность для фьючерсной цены, но входит в дисконтирование ожидаемой стоимости опциона.}

\subsubsection{Выплата на конечных узлах}

На экспирации стоимость опциона равна внутренней стоимости.

Для call:

\begin{equation}
V_N=\max(F_N-K,0).
\end{equation}

Здесь $V_N$ -- стоимость опциона в конечном узле дерева, то есть на экспирации. $F_N$ -- фьючерсная цена в этом конечном узле. Для call выплата положительна только тогда, когда $F_N>K$.

Для put:

\begin{equation}
V_N=\max(K-F_N,0).
\end{equation}

Для put всё наоборот: выплата положительна, когда фьючерсная цена ниже страйка.

\subsubsection{Обратная индукция}

После расчёта выплат на последнем слое дерево сворачивается назад. Для европейского опциона стоимость в узле равна дисконтированному риск-нейтральному ожиданию двух будущих значений:

\begin{equation}
V =
e^{-r\Delta t}
\left(qV^{\uparrow}+(1-q)V^{\downarrow}\right).
\end{equation}

Здесь $V^{\uparrow}$ -- стоимость опциона в следующем узле, если фьючерс пошёл вверх, а $V^{\downarrow}$ -- стоимость опциона в следующем узле, если фьючерс пошёл вниз. Выражение $qV^{\uparrow}+(1-q)V^{\downarrow}$ -- риск-нейтральное ожидаемое значение опциона на следующем шаге. Множитель $e^{-r\Delta t}$ переводит эту будущую ожидаемую стоимость в текущую стоимость.

Для американского опциона добавляется проверка раннего исполнения:

\begin{equation}
V =
\max\left(
e^{-r\Delta t}
\left(qV^{\uparrow}+(1-q)V^{\downarrow}\right),
H(F,K)
\right).
\end{equation}

Смысл: в каждом узле держатель выбирает лучшее из двух действий: продолжить держать опцион или исполнить его прямо сейчас. Поэтому цена американского опциона не может быть меньше внутренней стоимости.

Почему это называется обратной индукцией: мы не можем сразу узнать сегодняшнюю цену, пока не знаем, сколько опцион будет стоить завтра в каждом возможном состоянии. Но на самой экспирации цену знать легко: она равна payoff. Поэтому алгоритм начинает с конца, где всё известно, и шаг за шагом идёт назад к сегодняшнему дню.

Мини-пример. Пусть до экспирации один шаг. На последнем шаге call может стоить 20 в верхнем состоянии и 0 в нижнем. Если $q=0.5$, $r=0$ и досрочного исполнения нет, сегодняшняя цена равна $0.5\cdot20+0.5\cdot0=10$. Если шагов два, сначала считаются цены в двух предпоследних узлах по последним выплатам, а потом из этих двух значений считается цена в корне.

\answer{Если спросят, что такое continuation value: это стоимость «не исполнять сейчас, а держать опцион дальше». В CRR она равна дисконтированному риск-нейтральному ожиданию будущих стоимостей. Для американского опциона мы сравниваем continuation value с immediate exercise value.}

\subsubsection{Алгоритм CRR в работе}

Для каждого наблюдения:

\begin{enumerate}
    \item Проверяется корректность входов: $F>0$, $K>0$, $T>0$, $\sigma>0$, тип опциона call/put.
    \item Считается $\Delta t=T/N$.
    \item Считаются $u$, $d$ и $q$.
    \item Строятся фьючерсные цены на последнем слое дерева.
    \item Считаются terminal payoffs.
    \item Дерево сворачивается назад: continuation value дисконтируется, затем сравнивается с immediate exercise value.
    \item Значение в корне дерева является текущей модельной ценой.
\end{enumerate}

\impl{В \texttt{options\_utils.py} функции \texttt{black76\_price} и \texttt{american\_futures\_option\_crr} написаны вручную. Black-76 реализован векторно через \texttt{numpy} и функцию нормального распределения \texttt{scipy.stats.norm.cdf}. CRR-дерево также запрограммировано вручную: сначала строится последний слой дерева, затем цикл \texttt{for step in range(...)} сворачивает дерево назад. Стандартная библиотека \texttt{arch} используется только для оценки GARCH, а не для ценообразования опционов.}

\answer{Если спросят, почему CRR выбран основной моделью: она проста, интерпретируема, численно реализуема и умеет учитывать досрочное исполнение, что критично для американских опционов.}

\section{Четыре модельные спецификации}

\subsection{Три американские спецификации}

Основное сравнение проводится между:

\begin{itemize}
    \item CRR + HV\_21d;
    \item CRR + HV\_63d;
    \item CRR + GARCH.
\end{itemize}

Во всех трёх случаях модель цены одна и та же -- американская CRR. Меняется только входная волатильность $\sigma$.

Это важно методологически: так можно отделить эффект способа оценки волатильности от эффекта класса модели.

\subsection{Европейский benchmark}

Дополнительно рассчитывается:

\begin{itemize}
    \item Black-76 + HV\_63d.
\end{itemize}

Он сравнивается с CRR + HV\_63d, чтобы отличие было именно в стиле исполнения, а не в волатильности.

\begin{equation}
\Pi^R_i =
V^{\text{CRR}}_{63,i}
-
V^{\text{Black}}_{63,i}.
\end{equation}

$\Pi^R_i$ -- премия раннего исполнения для наблюдения $i$. $V^{\text{CRR}}_{63,i}$ -- цена американского опциона по CRR с 63-дневной исторической волатильностью, а $V^{\text{Black}}_{63,i}$ -- цена европейского опциона Black-76 с той же волатильностью. Если премия велика, значит право досрочного исполнения существенно влияет на стоимость опциона.

\section{Метрики ошибок}

\subsection{Знаковая ошибка}

Базовая ошибка:

\begin{equation}
e_i=\widehat{P}_i-P_i.
\end{equation}

Она сохраняет знак. Поэтому по ней можно понять bias:

\begin{itemize}
    \item $e_i<0$ -- модель недооценивает рыночную цену;
    \item $e_i>0$ -- модель переоценивает рыночную цену.
\end{itemize}

\subsection{MAE}

Средняя абсолютная ошибка:

\begin{equation}
MAE =
\frac{1}{n}
\sum_{i=1}^{n}
|e_i|.
\end{equation}

Здесь $n$ -- число наблюдений, а $|e_i|$ -- абсолютная ошибка по наблюдению $i$. MAE показывает средний масштаб ошибки в пунктах цены. Её легко интерпретировать: если MAE равно 107, то модель в среднем ошибается примерно на 107 пунктов.

Плюс MAE: устойчива к знаку, не даёт положительным и отрицательным ошибкам взаимно уничтожиться. Минус: не показывает направление bias.

\subsection{RMSE}

Среднеквадратичная ошибка:

\begin{equation}
RMSE =
\sqrt{
\frac{1}{n}
\sum_{i=1}^{n}
e_i^2
}.
\end{equation}

Здесь $e_i^2$ -- квадрат ошибки. RMSE сильнее штрафует крупные ошибки, потому что ошибка сначала возводится в квадрат. Если у модели есть редкие, но очень большие промахи, RMSE будет расти сильнее, чем MAE.

\answer{Если спросят, почему у GARCH лучший MAE, а у HV\_63d лучший RMSE: GARCH лучше по типичной абсолютной ошибке, но HV\_63d чуть устойчивее к крупным выбросам. Поэтому они лидируют по разным критериям.}

\subsection{Mean error}

Средняя знаковая ошибка:

\begin{equation}
\overline{e} =
\frac{1}{n}
\sum_{i=1}^{n}
e_i.
\end{equation}

Здесь ошибки сохраняют знак, поэтому положительные и отрицательные значения могут компенсироваться. Эта метрика отвечает на вопрос: есть ли систематическое смещение. Если $\overline{e}$ отрицательна, модель систематически занижает цену. Если положительна, завышает.

В работе у всех американских моделей mean error отрицателен. Это означает систематическую недооценку рынка.

\subsection{MAPE}

Средняя абсолютная процентная ошибка:

\begin{equation}
MAPE =
\frac{1}{n}
\sum_{i=1}^{n}
\left|
\frac{e_i}{P_i}
\right|.
\end{equation}

Здесь $\frac{e_i}{P_i}$ -- ошибка как доля от рыночной цены опциона. MAPE нормирует ошибку на полную рыночную цену. Это удобно, когда цены контрактов сильно различаются по масштабу. Но для опционов есть проблема: полная цена может состоять в основном из внутренней стоимости, которая не является тем компонентом, где волатильностная модель наиболее важна.

\section{TVNAE: авторская метрика по временной стоимости}

\subsection{Формула}

TVNAE определяется так:

\begin{equation}
TVNAE_i =
\frac{|\widehat{P}_i-P_i|}{TV_i},
\qquad
TV_i=P_i-H(F_i,K_i).
\end{equation}

В расчётах используется фильтр:

\begin{equation}
TV_i \geq 10.
\end{equation}

Он нужен, чтобы не делить на слишком маленькое число. Если временна́я стоимость близка к нулю, даже небольшая абсолютная ошибка даст огромный и нестабильный TVNAE.

\subsection{Почему MAE может маскировать качество модели}

Представим deep ITM call:

\begin{itemize}
    \item цена фьючерса сильно выше страйка;
    \item внутренняя стоимость очень велика;
    \item временна́я стоимость может быть относительно небольшой.
\end{itemize}

Если модель хорошо воспроизводит внутреннюю стоимость, но плохо оценивает временную, MAE может выглядеть приемлемо, потому что большая часть цены механически задана $F-K$. Но именно временна́я стоимость является компонентой, которую модель должна объяснять через $\sigma$, $T$ и $r$.

\key{TVNAE измеряет ошибку относительно той части цены, где модель действительно работает как модель риска и неопределённости. Поэтому она содержательно важнее MAE для deep ITM/OTM структуры выборки.}

\subsection{Как интерпретировать значения}

Если $TVNAE=0.66$, это означает, что абсолютная ошибка в среднем составляет 66\% от временной стоимости опциона. Если $TVNAE=0.19$, ошибка составляет 19\% от временной стоимости. Чем меньше TVNAE, тем лучше модель описывает опциональную компоненту цены.

\answer{Если спросят, почему TVNAE не заменяет все метрики: потому что это дополнительная диагностика. MAE/RMSE показывают ошибку в пунктах цены, а TVNAE показывает ошибку относительно временной стоимости. Они отвечают на разные вопросы.}

\section{Премия раннего исполнения}

\subsection{Формула}

Премия раннего исполнения рассчитывается как:

\begin{equation}
\Pi^R_i =
V^{\text{CRR}}_{63,i}
-
V^{\text{Black}}_{63,i}.
\end{equation}

Здесь обе модели используют одинаковую HV\_63d. Поэтому разница не связана с волатильностью. Она связана именно с тем, что CRR учитывает американское исполнение, а Black-76 -- нет.

\subsection{Экономический смысл}

Американский опцион даёт держателю больше прав, чем европейский: его можно исполнить раньше. Поэтому цена американского опциона должна быть не ниже цены аналогичного европейского опциона при прочих равных.

Для длинных опционов и высокой ставки право раннего исполнения особенно важно. Если опцион глубоко в деньгах, досрочное исполнение позволяет раньше получить внутреннюю стоимость и использовать высвобожденный капитал.

\subsection{Нижняя граница цены}

Американский опцион не должен стоить меньше внутренней стоимости:

\begin{equation}
V^{\text{American}} \geq H(F,K).
\end{equation}

Если Black-76 даёт цену ниже внутренней стоимости, это показывает, что европейская модель непригодна как основная модель для американского опциона. В работе такое происходит в 24.5\% наблюдений.

\answer{Если спросят, почему Black-76 ниже intrinsic -- это не ошибка кода, а свойство европейской модели: она не даёт права немедленного исполнения, поэтому может нарушать нижнюю границу, естественную для американского опциона.}

\section{Группировки наблюдений}

\subsection{Режимы волатильности}

Режимы строятся по 21-дневной исторической волатильности базовых фьючерсов. Распределение HV\_21d делится по терцилям:

\begin{itemize}
    \item нижняя треть -- low\_vol;
    \item средняя треть -- mid\_vol;
    \item верхняя треть -- high\_vol.
\end{itemize}

Такой подход не требует заранее фиксировать произвольные численные пороги. Режимы определяются эмпирически из распределения данных.

\subsection{DTE buckets}

Срок до экспирации делится на интервалы:

\begin{itemize}
    \item 0--7 дней;
    \item 8--30 дней;
    \item 31--90 дней;
    \item 91+ дней.
\end{itemize}

Это позволяет проверить, одинаково ли модели работают на коротких, средних и длинных опционах.

\subsection{Moneyness buckets}

Денежность делится на:

\begin{itemize}
    \item deep\_OTM: $F/K<0.9$;
    \item OTM: $0.9 \leq F/K <0.97$;
    \item ATM: $0.97 \leq F/K <1.03$;
    \item ITM: $1.03 \leq F/K <1.1$;
    \item deep\_ITM: $F/K \geq 1.1$.
\end{itemize}

Эта группировка особенно полезна, потому что ошибки моделей часто зависят от положения опциона относительно денег. ATM-опционы наиболее чувствительны к волатильности, а deep ITM опционы имеют большую внутреннюю стоимость.

\section{Тест средней ошибки}

\subsection{Что проверяется}

Тест средней ошибки проверяет, отличается ли средняя ошибка модели от нуля:

\begin{equation}
H_0:\ \mu=0,
\qquad
H_1:\ \mu\neq 0.
\end{equation}

Здесь $H_0$ -- нулевая гипотеза. Она говорит: в среднем модель не смещена, то есть иногда ошибается вверх, иногда вниз, но средняя ошибка равна нулю. $H_1$ -- альтернативная гипотеза: средняя ошибка не равна нулю, то есть есть систематический bias.

Если нулевая гипотеза не отвергается, то нет статистического доказательства систематического смещения. Если отвергается и $\mu<0$, модель статистически значимо недооценивает рынок.

\subsection{Почему выбран именно этот тест}

В работе важно не только сравнить абсолютный размер ошибки, но и понять, есть ли направленное смещение. MAE и RMSE всегда неотрицательны и не показывают, модель завышает или занижает цену. Поэтому нужен отдельный тест по знаковой ошибке $e_i$.

Если средняя ошибка отрицательна и статистически значима, это говорит о систематической недооценке рынка. Это сильнее, чем просто сказать «среднее отрицательное в таблице»: тест показывает, что такой результат маловероятно объясняется случайной вариацией выборки.

\subsection{Регрессионная форма теста}

В работе тест записан как регрессия ошибки на константу:

\begin{equation}
e_{i,t}=\mu+\varepsilon_{i,t}.
\end{equation}

Здесь $e_{i,t}$ -- ошибка $i$-го опциона в торговый день $t$, $\mu$ -- средняя ошибка, которую мы оцениваем, а $\varepsilon_{i,t}$ -- остаток, то есть индивидуальное отклонение конкретной ошибки от среднего уровня.

Оценка константы $\hat{\mu}$ равна выборочной средней ошибке. t-статистика:

\begin{equation}
t =
\frac{\hat{\mu}}{SE(\hat{\mu})}.
\end{equation}

Здесь $SE(\hat{\mu})$ -- стандартная ошибка оценки средней ошибки. Чем больше $|\hat{\mu}|$ относительно своей стандартной ошибки, тем больше по модулю t-статистика и тем сильнее свидетельство против $H_0$.

\subsection{Что такое p-value простыми словами}

p-value отвечает на вопрос: если бы нулевая гипотеза была верна, насколько вероятно получить настолько экстремальную t-статистику или ещё экстремальнее. Если p-value меньше 0.05, обычно говорят, что результат статистически значим на 5\% уровне. В работе p-value меньше 0.001, поэтому доказательство систематического смещения очень сильное.
\subsection{Почему стандартные ошибки кластеризуются по дате}

В один торговый день наблюдается много опционов. Их ошибки не независимы: на все контракты этого дня влияют общие рыночные условия, уровень фьючерса, волатильность и новости.

Если считать такие наблюдения полностью независимыми, стандартные ошибки могут быть занижены, а статистическая значимость завышена. Поэтому используются clustered standard errors по торговой дате.

\key{Кластеризация по дате говорит: внутри одного дня ошибки могут быть связаны, но между разными датами мы используем больше независимой вариации.}

\impl{В \texttt{options\_utils.py} это функция \texttt{mean\_error\_test}. Она строит регрессию \texttt{error \textasciitilde{} 1} через \texttt{statsmodels.formula.api.ols}. Стандартные ошибки задаются как \texttt{cov\_type="cluster"} с группами по \texttt{TRADEDATE}. То есть сам тест реализован стандартной OLS-функцией, но спецификация и кластеризация заданы вручную.}

\subsection{Как интерпретировать результат}

В работе все три американские модели имеют отрицательную среднюю ошибку и p-value меньше 0.001. Это означает, что систематическое занижение статистически значимо.

\answer{Если спросят, что именно доказывает этот тест: он доказывает не то, что модель «плохая вообще», а то, что её средняя ошибка статистически не равна нулю. В данном случае она отрицательная, значит есть систематическая недооценка.}

\section{Тесты Диболда--Мариано}

\subsection{Зачем нужен DM-test}

Тест средней ошибки говорит, есть ли bias у одной модели. Но он не отвечает на вопрос, значимо ли одна модель лучше другой по функции потерь. Для этого используются тесты Диболда--Мариано.

DM-test сравнивает прогнозную точность двух моделей. В работе сравниваются дневные средние функции потерь:

\begin{itemize}
    \item абсолютная ошибка -- для MAE-loss;
    \item квадратичная ошибка -- для MSE-loss.
\end{itemize}

\subsection{Почему выбран именно DM-test}

В этой работе есть несколько моделей, и простой факт «у одной MAE меньше» ещё не означает, что она статистически лучше. Ошибки меняются по дням, по контрактам и по рыночным режимам. DM-test выбран потому, что это стандартный тест для сравнения прогнозной точности двух моделей по выбранной функции потерь.

Он особенно подходит здесь, потому что модели оцениваются на одних и тех же наблюдениях. То есть для каждого дня можно сравнить, какая модель дала меньшую среднюю потерю, и затем проверить, является ли средняя разность потерь устойчивой.

\subsection{Разность потерь}

Для моделей $A$ и $B$ на каждый торговый день считается средняя потеря по всем контрактам этого дня:

\begin{equation}
\bar L^A_t,
\qquad
\bar L^B_t.
\end{equation}

Здесь $\bar L^A_t$ -- средняя потеря модели $A$ за день $t$, а $\bar L^B_t$ -- средняя потеря модели $B$ за тот же день. Потеря может быть абсолютной ошибкой $|e|$ или квадратичной ошибкой $e^2$.

Затем строится разность:

\begin{equation}
d_t=\bar L^A_t-\bar L^B_t.
\end{equation}

Если $d_t$ положительно, в этот день модель $A$ хуже, потому что её потеря больше. Если отрицательно, хуже модель $B$. Если вокруг нуля без устойчивого направления, модели статистически близки по качеству.

Если $d_t>0$, модель $A$ в этот день ошибалась сильнее, чем модель $B$. Если $d_t<0$, модель $A$ была лучше.

\subsection{Гипотеза теста}

Нулевая гипотеза:

\begin{equation}
H_0:\ \mathbb{E}[d_t]=0.
\end{equation}

Здесь $\mathbb{E}[d_t]$ -- математическое ожидание дневной разности потерь. Нулевая гипотеза говорит: в среднем модели имеют одинаковую потерю.

Это означает, что две модели имеют одинаковую ожидаемую функцию потерь.

Альтернатива:

\begin{equation}
H_1:\ \mathbb{E}[d_t]\neq 0.
\end{equation}

Альтернатива говорит: одна из моделей систематически имеет большую или меньшую потерю. В работе используется двусторонняя логика, потому что заранее формально проверяется не только «GARCH лучше», а вообще наличие различия.

\subsection{Регрессионная форма}

В работе тест реализован как регрессия:

\begin{equation}
d_t=\delta+u_t.
\end{equation}

Здесь $\delta$ -- средняя разность потерь, а $u_t$ -- отклонение дневной разности от среднего. Оценка $\hat{\delta}$ -- это средняя разность потерь. Если она статистически значима, модели различаются по точности.

\subsection{Почему HAC Newey--West}

Дневные разности потерь могут быть автокоррелированы: если сегодня модель ошибалась сильно из-за рыночного режима, завтра похожий режим может сохраниться. Также дисперсия разности потерь может быть непостоянной.

Поэтому используются HAC standard errors с 5 лагами. HAC означает heteroskedasticity and autocorrelation consistent: стандартные ошибки устойчивы к гетероскедастичности и автокорреляции.

Пять лагов означает, что при расчёте стандартной ошибки учитывается зависимость дневных разностей потерь примерно на несколько предыдущих торговых дней. Это важно, потому что ошибки моделей на финансовых данных часто имеют серийную зависимость.

\subsection{Как читать знак результата}

В таблице сравнение записывается как «model A minus model B». Если средняя разность положительна и значима, модель A хуже модели B. Если отрицательна и значима, модель A лучше модели B. Если p-value большое, статистически значимой разницы нет.

В работе:

\begin{itemize}
    \item HV\_21d значимо хуже HV\_63d по MAE;
    \item GARCH и HV\_63d статистически неразличимы на полной выборке;
    \item по MSE различие GARCH и HV\_21d слабее и не проходит 5\% уровень.
\end{itemize}

\answer{Если спросят, почему GARCH численно чуть лучше по MAE, но статистически неотличим от HV\_63d: потому что разница в средней потере мала относительно вариации дневных разностей потерь. Численное лидерство не всегда является статистически значимым.}

\impl{В \texttt{options\_utils.py} это функция \texttt{dm\_test\_daily}. Сначала данные группируются по \texttt{TRADEDATE}, затем для каждой даты считается средняя потеря двух моделей. После этого строится \texttt{loss\_diff} и оценивается регрессия \texttt{loss\_diff \textasciitilde{} 1} через \texttt{statsmodels}. HAC-ошибки задаются как \texttt{cov\_type="HAC"} и \texttt{maxlags=5}.}

\section{OLS-регрессия абсолютной ошибки}

\subsection{Зачем нужна регрессия}

Регрессия в работе не является моделью цены опциона. Она не пытается предсказывать рыночную цену. Её задача описательная: понять, какие характеристики наблюдения связаны с величиной ошибки.

Иными словами, после того как мы увидели, что модели ошибаются, регрессия отвечает на следующий вопрос: ошибка больше из-за режима высокой волатильности, из-за длинного срока до экспирации, из-за денежности или из-за типа опциона? Это помогает не просто ранжировать модели, а понять структуру их слабых мест.

Зависимая переменная:

\begin{equation}
|e_i|.
\end{equation}

Модуль ошибки берётся потому, что здесь важен размер промаха, а не направление. Если бы зависимой переменной была знаковая ошибка $e_i$, положительные и отрицательные ошибки могли бы взаимно компенсироваться. Абсолютная ошибка показывает, насколько далеко модель от рынка в пунктах.

В основной регрессии используется абсолютная ошибка модели CRR + HV\_63d. Это устойчивый baseline, который статистически неотличим от GARCH на полной выборке и имеет лучший RMSE.

\subsection{Финальная спецификация}

\begin{equation}
|e_i|
=
\beta_0
+ \beta_1\mathbf{1}[\text{high\_vol}]_i
+ \beta_2DTE_i
+ \beta_3\left|\ln(F_i/K_i)\right|
+ \beta_4\mathbf{1}[\text{put}]_i
+ \varepsilon_i.
\end{equation}

Здесь $\beta_0$ -- базовый уровень ошибки для базовой категории наблюдений: не high-vol, call-опцион, при нулевых значениях числовых регрессоров в формальном смысле. Коэффициенты $\beta_1,\ldots,\beta_4$ показывают, как меняется ожидаемая абсолютная ошибка при изменении соответствующего фактора при прочих равных. Остаток $\varepsilon_i$ -- часть ошибки, которую эта простая описательная регрессия не объясняет.

\subsection{Почему выбран OLS}

OLS выбран как простой и интерпретируемый способ описать связь между величиной ошибки и характеристиками опциона. Его преимущество для защиты: каждый коэффициент легко объяснить в пунктах ошибки. Например, коэффициент при DTE показывает, сколько пунктов ошибки добавляет один дополнительный день до экспирации при прочих равных.

Да, абсолютная ошибка не обязана быть идеально нормально распределена, поэтому акцент делается не на идеальной predictive-модели, а на описательной интерпретации и робастных стандартных ошибках, кластеризованных по дате.
\subsection{Смысл переменных}

\begin{itemize}
    \item $\mathbf{1}[\text{high\_vol}]$ -- индикатор высокой волатильности. Равен 1, если наблюдение относится к high\_vol regime.
    \item $DTE$ -- число календарных дней до экспирации. Проверяет, растёт ли ошибка с горизонтом.
    \item $\left|\ln(F/K)\right|$ -- расстояние от ATM-зоны. Проверяет, как ошибка связана с денежностью.
    \item $\mathbf{1}[\text{put}]$ -- индикатор put-опциона. Базовая категория -- call.
\end{itemize}

Индикаторная переменная -- это переменная, принимающая значения 0 или 1. Если наблюдение не относится к high\_vol, индикатор равен 0 и вклад $\beta_1$ исчезает. Если относится, к прогнозируемой абсолютной ошибке добавляется $\beta_1$.

\subsection{Как интерпретировать коэффициенты}

Коэффициент $\beta_2$ при DTE показывает, на сколько пунктов в среднем меняется абсолютная ошибка при увеличении срока до экспирации на один день, при прочих равных.

В финальной спецификации коэффициент около $+0.16$. Значит, один дополнительный день до экспирации связан примерно с ростом абсолютной ошибки на 0.16 пункта.

Коэффициент при $|\ln(F/K)|$ отрицательный. Это означает, что после контроля DTE, режима и типа опциона ошибки больше около ATM и меньше дальше от ATM. На первый взгляд это может быть неочевидно, потому что deep ITM контракты имеют крупные абсолютные ошибки в простых таблицах. Но регрессия контролирует другие факторы, прежде всего DTE и тип опциона.

Коэффициент при put отрицательный. Это значит, что put-опционы в среднем оцениваются точнее call-опционов в этой спецификации.

Коэффициент при high\_vol после добавления DTE и денежности становится отрицательным. Это показывает, что сам режим высокой волатильности не является самостоятельной причиной роста ошибки. В простом разрезе high\_vol выглядит хуже из-за состава наблюдений: там больше длинных контрактов.

\warn{Нельзя говорить, что OLS доказывает причинность. Это описательная регрессия: она показывает статистическую связь характеристик с величиной ошибки.}

\subsection{Как читать t-stat и p-value в OLS}

Для каждого коэффициента проверяется гипотеза, что он равен нулю:

\begin{equation}
H_0:\ \beta_j=0.
\end{equation}

Если p-value маленькое, коэффициент статистически отличен от нуля. Это означает, что соответствующий фактор связан с абсолютной ошибкой даже после контроля остальных факторов в регрессии.

t-статистика считается как:

\begin{equation}
t_j=\frac{\hat{\beta}_j}{SE(\hat{\beta}_j)}.
\end{equation}

Здесь $\hat{\beta}_j$ -- оценённый коэффициент, а $SE(\hat{\beta}_j)$ -- его стандартная ошибка. Чем больше t-статистика по модулю, тем сильнее свидетельство, что коэффициент не равен нулю.

\subsection{Почему используются clustered SE}

Как и в тесте средней ошибки, стандартные ошибки кластеризуются по торговой дате. Причина та же: ошибки опционов внутри одного дня могут быть связаны общими рыночными шоками. Кластеризация делает выводы о значимости более корректными.

\subsection{Что означает R-squared 0.267}

$R^2$ показывает долю вариации зависимой переменной, объясняемую регрессорами. Значение 0.267 означает, что модель объясняет около 27\% вариации абсолютной ошибки.

Для финансовых микроданных это содержательный результат: ошибка опционной модели зависит от множества ненаблюдаемых факторов, поэтому ожидать $R^2$ близкий к единице не нужно.

\answer{Если спросят, почему $R^2$ не очень высокий: потому что регрессия описывает ошибку, а не детерминированную цену. На ошибку влияют ликвидность, микроструктура, рыночные ожидания, спрос/предложение и другие факторы, которых нет в модели.}

\impl{В \texttt{calculations.ipynb} строятся три OLS-спецификации. Первая включает только \texttt{high\_vol\_dummy}. Вторая добавляет \texttt{DTE\_DAYS} и \texttt{abs\_log\_moneyness}. Третья добавляет тип опциона через \texttt{C(option\_type\_norm)}. Все модели оцениваются через \texttt{statsmodels.formula.api.ols}, а стандартные ошибки задаются как \texttt{cov\_type="cluster"} с группировкой по торговой дате.}

\section{Проверки устойчивости}

\subsection{Зачем нужны robustness checks}

Проверки устойчивости отвечают на вопрос: сохраняется ли вывод, если изменить состав выборки или параметр расчёта. Если результат держится только на одной специфической подвыборке, он слабее. Если вывод сохраняется в разных условиях, он надёжнее.

\subsection{Подвыборки}

В работе проверяются:

\begin{itemize}
    \item полная выборка;
    \item выборка без deep OTM и deep ITM;
    \item короткие и среднесрочные опционы с DTE $\leq 91$;
    \item длинные опционы с DTE $>91$;
    \item выборка без high\_vol режима.
\end{itemize}

Основной вывод: GARCH особенно хорош на коротких и среднесрочных контрактах, а на длинных HV\_63d и GARCH становятся почти равны. Это согласуется с экономической логикой: GARCH даёт краткосрочный условный прогноз, а длинные опционы требуют информации о долгом горизонте волатильности.

\subsection{Чувствительность к ставке}

Базовая ставка равна 16\%. Дополнительно проверяются 12\% и 20\%. Это нужно, потому что ставка влияет на дисконтирование и на ценность раннего исполнения.

Качественное ранжирование моделей при изменении ставки сохраняется. Значит, выводы не являются артефактом одного выбранного значения $r$.

\answer{Если спросят, почему ставка фиксированная: это упрощение базовой модели. В будущей работе лучше использовать срочную структуру ставок, но robustness показывает, что основной вывод не ломается при разумном изменении ставки.}

\section{Как объяснять основные результаты методологически}

\subsection{Почему все модели недооценивают рынок}

Главное объяснение -- слишком простая структура волатильности. В моделях используется одна $\sigma$ на контракт, но реальный рынок имеет поверхность подразумеваемой волатильности:

\begin{equation}
\sigma_{\text{impl}}=f(K,T).
\end{equation}

Одна плоская волатильность не может одновременно правильно оценить опционы разных страйков и сроков. Кроме того, рыночные цены могут включать ликвидностную премию, премию за хвостовые риски и эффекты спроса на защиту.

\subsection{Почему HV\_21d хуже}

HV\_21d быстро реагирует на рынок, но она шумная. В выборке много длинных опционов, а для них краткосрочное окно не отражает весь горизонт неопределённости до экспирации. Поэтому HV\_21d статистически хуже HV\_63d на полной выборке.

\subsection{Почему GARCH выигрывает на коротких DTE}

GARCH оценивает текущую условную волатильность и хорошо улавливает недавнюю турбулентность. Для коротких опционов именно текущий режим особенно важен, потому что до экспирации осталось мало времени. Поэтому GARCH резко лучше по TVNAE на DTE 0--7 и 8--30.

\subsection{Почему GARCH не выигрывает в high\_vol}

В high\_vol режиме GARCH может быть инерционным. Если шок резкий, модель может либо медленно адаптироваться, либо слишком долго переносить текущую турбулентность в прогноз. Кроме того, high\_vol сегмент в выборке связан с длинными контрактами, где однодневный GARCH-прогноз менее релевантен.

\subsection{Почему класс модели важнее оценки волатильности}

Сравнение CRR + HV\_63d и Black-76 + HV\_63d изолирует эффект early exercise. Волатильность одинаковая, но CRR учитывает досрочное исполнение, а Black-76 нет. Разрыв по MAE значительно больше, чем разница между американскими волатильностными спецификациями.

\key{Иерархический вывод: сначала нужно выбрать правильный класс модели для американского опциона, и только потом оптимизировать способ оценки волатильности.}

\section{Как именно это реализовано в коде}

\subsection{Разделение ролей между файлами}

Практическая часть находится в двух основных файлах:

\begin{itemize}
    \item \texttt{options\_utils.py} -- файл с функциями: Black-76, CRR, GARCH forecast, расчёт ошибок, метрик и статистических тестов;
    \item \texttt{calculations.ipynb} -- ноутбук, где эти функции применяются к данным, строятся выборки, таблицы, графики и сохраняются результаты.
\end{itemize}

То есть \texttt{options\_utils.py} -- это «инструменты», а \texttt{calculations.ipynb} -- «сценарий исследования».

\subsection{Что написано вручную}

Вручную запрограммированы:

\begin{itemize}
    \item формула Black-76 в функции \texttt{black76\_price};
    \item американское CRR-дерево в функции \texttt{american\_futures\_option\_crr};
    \item расчёт модельных ошибок в \texttt{add\_model\_block};
    \item отбор общей сопоставимой выборки в \texttt{restrict\_to\_common\_sample};
    \item расчёт MAE, RMSE, mean error, MAPE в \texttt{calc\_metrics};
    \item расчёт TVNAE в \texttt{compute\_tvnae\_series} и \texttt{calc\_tvnae\_metrics};
    \item логика expanding-window прогноза GARCH вокруг стандартной оценки модели;
    \item группировки по DTE, moneyness и vol regime в ноутбуке.
\end{itemize}

\subsection{Где использовались стандартные библиотеки}

Использовались стандартные Python-библиотеки:

\begin{itemize}
    \item \texttt{numpy} -- векторные вычисления, экспоненты, логарифмы, корни, массивы;
    \item \texttt{pandas} -- таблицы, фильтрация, группировки, merge, rolling windows;
    \item \texttt{scipy.stats.norm.cdf} -- функция распределения стандартного нормального закона в Black-76;
    \item \texttt{arch.arch\_model} -- оценка параметров GARCH(1,1);
    \item \texttt{statsmodels.formula.api.ols} -- OLS-регрессии, mean error test и DM-test в регрессионной форме;
    \item \texttt{matplotlib} -- графики.
\end{itemize}

\answer{Если спросят, «модели были из библиотеки или написаны вручную»: Black-76 и CRR-ценообразование написаны вручную. GARCH оценивался через стандартную библиотеку \texttt{arch}, но логика rolling/expanding forecast была прописана вручную. Статистические тесты были реализованы через стандартный OLS-интерфейс \texttt{statsmodels} с нужными типами робастных стандартных ошибок.}

\subsection{Как идёт расчёт в ноутбуке}

В \texttt{calculations.ipynb} сначала формируется \texttt{modelling\_df}: очищенная таблица наблюдений с рыночными ценами, фьючерсными ценами, страйками, сроками и типами опционов. Затем строится \texttt{underlying\_daily}: ряд базовых фьючерсов по датам, на котором считаются \texttt{ret\_1d}, \texttt{hv\_21d}, \texttt{hv\_63d} и режимы волатильности.

После этого вызывается \texttt{compute\_expanding\_garch\_forecast}, и прогноз GARCH присоединяется обратно к таблице наблюдений. Затем функция \texttt{add\_model\_block} по очереди добавляет четыре набора модельных цен и ошибок:

\begin{itemize}
    \item \texttt{crr\_hv\_21d};
    \item \texttt{crr\_hv\_63d};
    \item \texttt{crr\_garch};
    \item \texttt{black76\_hv\_63d}.
\end{itemize}

Дальше создаются группы \texttt{DTE\_bucket}, \texttt{moneyness\_bucket} и общие выборки \texttt{american\_common} и \texttt{all\_models\_common}. На них считаются итоговые таблицы, TVNAE, премия раннего исполнения, тесты и регрессии.

\subsection{Почему такая реализация защищает от look-ahead bias}

Look-ahead bias возник бы, если бы для цены опциона в дату $t$ использовались будущие данные после даты $t$. В работе это предотвращается так:

\begin{itemize}
    \item историческая волатильность считается rolling window по прошлым доходностям;
    \item GARCH оценивается на expanding window из уже накопленных наблюдений;
    \item прогноз GARCH на дату строится только после накопления прошлой истории;
    \item параметры GARCH переоцениваются периодически, но без использования будущих доходностей.
\end{itemize}

Это важно для честного сравнения моделей: каждая модель должна использовать только ту информацию, которая была бы доступна на момент оценки.

\section{Короткие ответы на вероятные вопросы}

\subsection*{Почему не использовалась подразумеваемая волатильность?}

Подразумеваемая волатильность извлекается из рыночной цены опциона. Если использовать её как вход в модель и потом сравнивать модельную цену с той же рыночной ценой, возникнет частичная тавтология. В работе цель -- проверить, как модели с наблюдаемыми историческими/GARCH оценками воспроизводят рынок, а не откалибровать модель точно под рыночную цену.

\subsection*{Что такое риск-нейтральная вероятность одним предложением?}

Это не реальная вероятность события, а такие модельные веса, при которых цена производного инструмента получается без арбитража: будущие выплаты усредняются так, будто инвесторы не требуют премию за риск, а затем дисконтируются по безрисковой ставке.

\subsection*{Почему в формуле Black--Scholes исчезает ожидаемая доходность актива?}

Потому что в выводе строится дельта-хеджированный портфель, где случайная часть риска локально убрана. Такой портфель становится безрисковым и должен зарабатывать ставку $r$. После хеджирования параметр $\mu$ сокращается, и цена опциона зависит от $r$, $\sigma$, $S$, $K$ и $T$, но не от субъективной ожидаемой доходности.

\subsection*{Что такое обратная индукция совсем просто?}

Это расчёт «с конца назад». На экспирации выплата опциона известна точно. Поэтому сначала считаются последние значения, потом предпоследние как дисконтированное ожидание последних, затем предыдущие, и так до сегодняшнего узла.

\subsection*{Почему биномиальное дерево можно считать нормальной аппроксимацией?}

Потому что много маленьких движений вверх/вниз в лог-ценах при большом числе шагов начинает приближать нормальное распределение лог-доходности. Это дискретный аналог непрерывной модели, а при увеличении числа шагов CRR сходится к Black--Scholes для европейского случая.

\subsection*{Почему в CRR для фьючерса вероятность $q=(1-d)/(u-d)$, а не формула с $e^{r\Delta t}$?}

Потому что под риск-нейтральной мерой фьючерсная цена является мартингалом: ожидаемая фьючерсная цена на следующем шаге равна текущей. Поэтому условие записывается как $F=qFu+(1-q)Fd$, и после деления на $F$ получается $q=(1-d)/(u-d)$.

\subsection*{Почему рыночная цена -- расчётная цена биржи?}

Расчётная цена доступна систематически и используется как воспроизводимый бенчмарк. Последняя сделка или котировка могут быть шумными, особенно для менее ликвидных опционов. Но это ограничение: расчётная цена не является абсолютной фундаментальной стоимостью.

\subsection*{Почему CRR использует 100 шагов?}

100 шагов -- компромисс между точностью и вычислительной нагрузкой. В выборке десятки тысяч наблюдений, и для каждого нужно построить дерево. Увеличение числа шагов может повысить точность аппроксимации, но увеличит время расчёта. Для сравнительного исследования 100 шагов даёт разумный баланс.

\subsection*{Почему GARCH оценивался через библиотеку, а CRR вручную?}

GARCH -- статистическая модель с численной оптимизацией параметров, поэтому надёжнее использовать стандартную библиотеку \texttt{arch}. CRR и Black-76 являются прямыми формулами/алгоритмами ценообразования, поэтому их удобно и прозрачно реализовать вручную, контролируя именно фьючерсную риск-нейтральную вероятность и проверку досрочного исполнения.

\subsection*{Что означает negative mean error?}

Это значит, что модельная цена в среднем ниже рыночной:

\begin{equation}
\widehat{P}_i-P_i<0.
\end{equation}

Если это статистически значимо, модель систематически недооценивает рынок.

\subsection*{Почему TVNAE считается только при $TV\geq10$?}

Чтобы избежать нестабильности деления на почти ноль. Если временна́я стоимость равна 1 пункту, ошибка 5 пунктов даст TVNAE 5, хотя в абсолютном смысле ошибка мала. Порог делает метрику устойчивее.

\subsection*{Что показывает DM-test простыми словами?}

Он отвечает на вопрос: является ли разница в ошибках двух моделей систематической или могла возникнуть случайно из-за вариации по дням. Это тест равенства ожидаемых функций потерь.

\subsection*{Почему DM-test считался на дневных средних потерях, а не на каждой строке?}

Потому что внутри одного дня много опционов, и их ошибки зависят от общего рыночного состояния. Дневное агрегирование делает единицей сравнения торговый день и снижает проблему внутридневной зависимости наблюдений.

\subsection*{Почему для DM-test использованы HAC-ошибки?}

Потому что дневные разности потерь могут быть зависимы во времени: если модель плохо работает в определённом режиме сегодня, похожая ситуация может сохраниться несколько дней. HAC Newey--West учитывает автокорреляцию и гетероскедастичность в таких дневных рядах.

\subsection*{Что показывает OLS-регрессия простыми словами?}

Она показывает, какие характеристики опциона связаны с большой ошибкой модели. В работе главный фактор -- срок до экспирации, а не сам по себе high-vol regime.

\subsection*{Почему регрессия не доказывает причинность?}

Потому что это не эксперимент и не структурная модель цены. Регрессия показывает условные статистические связи в наблюдаемых данных. Например, положительный коэффициент при DTE означает, что длинные сроки связаны с большей ошибкой при прочих равных, но не доказывает единственный причинный механизм.

\subsection*{Почему ставка влияет на early exercise?}

Право досрочного исполнения связано со стоимостью денег во времени. Если опцион глубоко в деньгах, немедленное исполнение позволяет раньше получить внутреннюю стоимость или высвободить капитал. При высокой ставке выгода от более раннего получения/использования денег выше, поэтому премия раннего исполнения может расти.

\subsection*{Почему все модели систематически недооценивают рынок?}

Главная причина -- слишком простая структура: одна волатильность на контракт вместо поверхности подразумеваемой волатильности. Дополнительно рыночная расчётная цена может включать ликвидностную премию, премию за хвостовые риски и эффекты спроса/предложения, которых нет в базовых моделях.

\subsection*{Главный методологический вывод}

Для американских опционов на фьючерс MXI нужно учитывать досрочное исполнение. Среди американских моделей GARCH полезен для коротких и среднесрочных контрактов и лучше описывает временну́ю стоимость, но на полной выборке статистически близок к HV\_63d. Все простые модели с одной $\sigma$ систематически недооценивают рынок, поэтому дальнейшее развитие должно идти к поверхности подразумеваемой волатильности, учёту ликвидности, стохастической волатильности и скачков.

\end{document}
